\ifdefined\ishandout
\documentclass[11pt,handout]{beamer}
\else
\documentclass[11pt]{beamer}
\fi

% 日本語対応
\usepackage{luatexja}
\usepackage{luatexja-fontspec}
\renewcommand{\kanjifamilydefault}{\gtdefault}

\usepackage{mathptmx}
\renewcommand{\sfdefault}{lmss}
\renewcommand{\familydefault}{\sfdefault}
\usepackage[T1]{fontenc}
\usepackage{amsmath}
\usepackage{amssymb}
\usepackage{graphicx}
\PassOptionsToPackage{normalem}{ulem}
\usepackage{ulem}
\usepackage{caption}
\captionsetup{labelformat=empty}
\usepackage{bbm}
\usepackage{upgreek}
\setbeamertemplate{section in toc}[sections numbered]
\makeatletter
\usepackage{caption} 
\captionsetup[table]{skip=10pt}
\newcommand\makebeamertitle{\frame{\maketitle}}%
\AtBeginDocument{%
   \let\origtableofcontents=\tableofcontents
   \def\tableofcontents{\@ifnextchar[{\origtableofcontents}{\gobbletableofcontents}}
   \def\gobbletableofcontents#1{\origtableofcontents}
}

\def\Tiny{\fontsize{7pt}{8pt}\selectfont}
\def\Normal{\fontsize{8pt}{10pt}\selectfont}

\usetheme{Madrid}
\usecolortheme{lily}
\useinnertheme{rounded}

\setbeamertemplate{footline}{\hfill\Normal{\insertframenumber/\inserttotalframenumber}}
\setbeamertemplate{navigation symbols}{}

\newenvironment{changemargin}[2]{%
\begin{list}{}{%
\setlength{\topsep}{0pt}%
\setlength{\leftmargin}{#1}%
\setlength{\rightmargin}{#2}%
\setlength{\listparindent}{\parindent}%
\setlength{\itemindent}{\parindent}%
\setlength{\parsep}{\parskip}% 
}%
\item[]}{\end{list}}

\setbeamertemplate{footline}{\hfill\insertframenumber/\inserttotalframenumber}
\setbeamertemplate{navigation symbols}{}

\usepackage{graphicx}
\usepackage{epsfig}
\usepackage{bm}
\usepackage{epsf}
\usepackage{float}
\usepackage[final]{pdfpages}
\usepackage{multirow}
\usepackage{colortbl}
\usepackage{xkeyval}
\usepackage{listings}
\usepackage{ifthen}
\usepackage{tikz}

\usepackage{setspace}
\usepackage{ragged2e}

\setbeamersize{text margin left=1em,text margin right=1em} 

% Table formatting
\usepackage{booktabs}

% Decimal align
\usepackage{dcolumn}
\newcolumntype{d}[0]{D{.}{.}{5}}

\global\long\def\expec#1{\mathbb{E}\left[#1\right]}
\global\long\def\var#1{\mathrm{Var}\left[#1\right]}
\global\long\def\cov#1{\mathrm{Cov}\left[#1\right]}
\global\long\def\prob#1{\mathrm{Prob}\left[#1\right]}
\global\long\def\one{\mathbf{1}}
\global\long\def\diag{\operatorname{diag}}
\global\long\def\expe#1#2{\mathbb{E}_{#1}\left[#2\right]}
\DeclareMathOperator*{\plim}{\text{plim}}

\usepackage{appendixnumberbeamer}
\renewcommand{\thefootnote}{}

\setbeamertemplate{footline}
        {
      \leavevmode%
\hfill
        \insertframenumber{}\hspace*{2ex}\vspace{1ex}
}

\definecolor{blue}{RGB}{0, 0, 210}
\definecolor{red}{RGB}{170, 0, 0}

\makeatother

\usepackage[english]{babel}

\usepackage{tikz}
\newcommand*\circled[1]{\tikz[baseline=(char.base)]{             \node[circle,ball color=structure.fg, shade,   color=white,inner sep=1.2pt] (char) {\tiny #1};}} 

\makeatletter
\let\save@measuring@true\measuring@true
\def\measuring@true{%
  \save@measuring@true
  \def\beamer@sortzero##1{\beamer@ifnextcharospec{\beamer@sortzeroread{##1}}{}}%
  \def\beamer@sortzeroread##1<##2>{}%
  \def\beamer@finalnospec{}%
}
\makeatother

\setbeamersize{text margin left= .8em,text margin right=1em} 
\newenvironment{wideitemize}{\itemize\addtolength{\itemsep}{10pt}}{\enditemize}
\newenvironment{wideitemizeshort}{\itemize}{\enditemize}

 \newcommand{\indep}{\perp\!\!\!\!\perp} 

\begin{document}

%% タイトルスライド
\begin{frame}[noframenumbering]{}
\vspace{0.5cm}
\title[]{第2章：確率と統計}
\author{Jonathan Roth}
\date{数理計量経済学 I \\ ブラウン大学 \\ } 
\titlepage {\small{}\ }\thispagestyle{empty} \vspace{-30pt}

\end{frame}
 
\begin{frame}{連絡事項}

\begin{itemize}
\item
課題（Problem Set）1を公開しました。締め切りは9月19日（金）午後4時、Gradescopeでの提出です。
\bigskip

\item
今週からTAセッションとオフィスアワー（OH）が始まります。時間と暫定的な場所についてはCanvasを確認してください。
\bigskip

\item
運営面で質問はありますか？
\end{itemize}	
\end{frame}

\begin{frame}{「全体像（Big Picture）」の復習}
\vspace{0.1cm}
役割分担を思い出しましょう：
\begin{center}
\includegraphics[width=0.9\linewidth]{BigPicture.jpg}
\end{center}
\vspace{-0.2cm}

\begin{itemize}
\item
\textbf{統計（Statistics）}：観察されたサンプルデータは、関心のある母集団の観察可能な特徴とどう関連しているか？
\smallskip

\item
\textbf{識別（Identification）}：母集団の観察可能な特徴は、関心のあるターゲット・パラメータとどう関連しているか？
\smallskip
\pause{}

\item
これら両方のタスクにおいて、データがどのように生成されるかを議論するための数学的な言語が必要です。それが\textbf{確率と統計}です。

\end{itemize}

\end{frame}

\begin{frame}
	\centering
	\includegraphics[width = 0.5 \linewidth]{poppy_stick.jpg}
\end{frame}

\begin{frame}{アウトライン}
	
	1. 確率変数と確率分布
	\vspace{0.8cm}
	
	2. 平均と分散
	\vspace{0.8cm}

	3. 実験における識別
	\vspace{0.8cm}
		
	4. 無作為抽出と標本平均
	
	\vspace{0.8cm}
	5. 仮説検定と推論
	
\end{frame}

\begin{frame}{確率変数}
\begin{wideitemize} 
\item
確率論は、\textbf{ランダムな過程（random processes）}の研究を定式化したものです。\pause{}
\bigskip

\item ランダムな過程の例には何があるでしょうか？\pause{}\smallskip
\begin{itemize}
	\item コイン投げ -- 表か裏か？ \smallskip
	\item 米国の世帯を無作為に調査 -- 所得はいくらか？ 
\end{itemize}
\pause{}
\bigskip

\item ランダムな過程の結果（実現値）を\textbf{確率変数（random variable）}と呼びます。

\end{wideitemize}

\end{frame}

\begin{frame}{用語の整理}

\begin{itemize}
\item \textbf{結果（Outcomes）}とは、ランダムな過程における互いに排他的な結果のことです（例：1回のコイン投げの結果は「表」と「裏」）。
\vspace{0.1cm}\pause{}
\item \textbf{確率（Probability）}とは、ある結果が起こる尤もしさを捉えたものです（すなわち、過程を繰り返したときの発生頻度）。
\vspace{0.1cm}\pause{}
\item \textbf{標本空間（Sample space）}とは、起こりうるすべての結果の集合です。
\vspace{0.1cm}\pause{}
\item \textbf{事象（Event）}とは、標本空間の部分集合です。その確率は、含まれる結果の確率の和になります。
\end{itemize}
\vspace{0.4cm}
\pause{}
例：2枚の公正なコインを投げます。\pause{}

\begin{itemize}
\item 起こりうる結果（標本空間）は何ですか？
\vspace{0.1cm}\pause{}
\item 少なくとも1枚が表になる（事象）確率はいくらですか？
\end{itemize}

\end{frame}

\begin{frame}{確率変数と累積分布関数 (CDF)}
\begin{wideitemize}
\item
\textbf{確率変数}とは、ランダムな過程を数値で要約したものです（形式的には、標本空間上で定義された実数値関数）。

\begin{itemize}
\item 例：$X$ を表が出た回数とする。
\end{itemize}
\pause{}

\item 実数値の確率変数は、その\textbf{累積分布関数（CDF）}によって特徴付けられます。

$$F(x) = Pr( X \leq x ) $$

\noindent これは、$X$ がある値 $x$ 以下になる確率を示します。
\medskip\pause{}
\begin{itemize}
\item 注意：通常、「$X$」のような確率変数の実現値（すなわちランダムではない数値）を表すには、「$x$」のような小文字を使用します。
\end{itemize}

\end{wideitemize}
\end{frame} 

\begin{frame}
	\centering
	\includegraphics[width = 0.8\linewidth]{cdf-income}
	
	\begin{itemize}
		\item
		2016年、米国の世帯の約半分は所得が5.6万ドル以下でした。\smallskip
\pause{}
		\item 形式的には：$F(56,516)=0.5$
	\end{itemize}
\end{frame}

\begin{frame}{サポート（Support）}
	
\begin{wideitemize}
	\item
	確率変数 $X$ の\textbf{サポート}（$\mathbb{X}$ と表記）とは、$X$ が取りうる値の集合です。
	
		\begin{itemize}
			\item 
			$X$ が1年間のうち雇用されていた月数なら、$\mathbb{X} = \{0,1,...,12\}$ です。
			
			\item
			$X$ が所得なら、$\mathbb{X} = \mathbb{R}_{\geq 0}$（近似的に）です。
		\end{itemize}
	\smallskip\pause{}

	\item
	$X$ のサポートが有限（例：$\{0,1\}$）なら、その $X$ は\textbf{離散型}であると言います。
	
	\item
	$X$ のサポートが連続体（例：$\mathbb{R}$ や $[0,1]$）なら、その $X$ は\textbf{連続型}の分布に従うと言います（厳密にはCDFが微分可能な場合）。
\end{wideitemize}

\end{frame}

\begin{frame}{密度関数と質量関数}

\begin{itemize}
	\item 
	$X$ が離散型の場合、サポートの各値を取る確率として\textbf{確率質量関数（PMF）}を定義します：
	\vspace{-0.2cm}
	$$p(x) = Pr(X =x )$$
	\pause
\vspace{-0.4cm}
	\begin{itemize}
	\item
	離散型確率変数のCDFは次のようになります：
	$$F(x) = \sum_{x' \leq x} p(x') $$ 
	\end{itemize}

	\pause
	\item
	連続型確率変数の場合、\textbf{確率密度関数（PDF）}を $f(x) = \dfrac{d}{dx} F(x)$ と定義します。これはCDFが次であることを意味します：
	$$F(x) = \int_{-\infty}^{x} f(t) dt$$
	
	\pause
	\item
	表記上の注意：$p(x)$ と $f(x)$ の両方がPDF/PMFとして使われます。
\end{itemize}

\end{frame}

\begin{frame}{}
\vspace{0.2cm}
離散型確率変数の例：Wi-Fi接続の失敗回数

\begin{itemize}
\item 以下はPMFです。CDFはどうなりますか？
\end{itemize}

\begin{center}
\includegraphics[scale=0.35]{simple_pmf.png}
\end{center}

\end{frame}

\begin{frame}{}
\vspace{0.2cm}
連続型確率変数の例：通勤時間

\begin{itemize}
\item 以下はPDFです。CDFはどうなりますか？
\end{itemize}

\begin{center}
\includegraphics[scale=0.35]{simple_pdf.png}
\end{center}

\end{frame}

\begin{frame}{PDF/CDF の性質}
\vspace{0.2cm}
\begin{itemize}
\item CDF $F(x)=Pr(X\le x)$ の主要な性質：
\smallskip

\begin{itemize}
\item 非減少：$x > x^\prime$ ならば $F(x)\ge F(x^\prime)$
\vspace{0.1cm}
\item $\lim_{x\rightarrow -\infty} F(x)=0$ および $\lim_{x\rightarrow +\infty}F(x)=1$ を満たす
\end{itemize}
\vspace{0.5cm}
\pause{}

\item PDF $f(x)=\frac{\partial}{\partial x}F(x)$ の対応する性質：
\smallskip

\begin{itemize}
\item 非負：すべての $x\in\mathbb{X}$ に対して $f(x)\ge 0$
\vspace{0.1cm}
\item $\int_{x\in\mathbb{X}}f(x)dx=1$ を満たす
\vspace{0.1cm}
\item PMF の場合：$\sum_{x\in\mathbb{X}}p(x)=1$
\end{itemize}
\end{itemize}

\end{frame}

\begin{frame}{ベルヌーイ分布}
\vspace{0.1cm}
重要な離散分布：ベルヌーイ分布 $X\in\{0,1\}$\pause{}

\begin{itemize}
\item 例：大学卒業の有無、コインが「表」になるかどうか（「ダミー変数」とも呼ばれます）
\vspace{0.1cm}\pause{}
\item PMF：
\vspace{-0.4cm}
\begin{align*}
p(x)=
\begin{cases}
1-\pi, & x=0\\
\pi, & x=1
\end{cases}
\end{align*}
（ただし $\pi\in[0,1]$）
\vspace{0.1cm}\pause
\item $\pi$ は $X$ の平均（期待値）です。これについては後ほど正式に定義します。
\vspace{0.1cm}\pause{}
\item $X\sim Bernoulli(\pi)$ と表記されます。
\vspace{0.1cm}\pause{}
\item 2値変数 $X$ の分布は、必ずこの形式になります。
\end{itemize}

\end{frame}

\begin{frame}{正規分布と一様分布}
\vspace{0.1cm}
重要な連続分布：正規分布 $X\in\mathbb{R}$

\begin{itemize}
\item 例：年収の対数（近似的に）
\vspace{0.1cm}
\item PDF：$f(x)=\frac{1}{\sqrt{2\pi\sigma^2}}\exp\left(-\frac{1}{2}\frac{(x-\mu)^2}{\sigma^2}\right)$ （ただし $x\in\mathbb{R}, \sigma>0$）
\vspace{0.1cm}
\item ここで $\mu$ は $X$ の平均、$\sigma^2$ は分散です。
\vspace{0.1cm}
\item $X\sim\mathrm{N}(\mu,\sigma^2)$ と表記されます。
\vspace{0.1cm}
\item 有用な性質：$X$ が正規分布に従うなら、任意の定数 $a, b$ に対して $aX+b$ も正規分布に従います。

	\centering
\includegraphics<1>[width = 0.4 \linewidth]{normalpdf.jpeg}

\pause
\end{itemize}
\vspace{0.3cm}
知っておくべき別の連続分布：一様分布 $X\in (a,b)$
\pause{}

\begin{itemize}
\item 例：無作為な宝くじ番号（近似的に）
\vspace{0.1cm}\pause{}
\item PDF：$f(x)=\frac{1}{b-a}$ （ただし $x\in(a,b)$）
\vspace{0.1cm}\pause{}
\item $X\sim\mathrm{U}(a,b)$ と表記されます。これも線形変換について閉じています。
\end{itemize}

\end{frame}

\begin{frame}{同時分布}
\vspace{0.2cm}
経済学の興味深い問いの多くは、2つ以上の確率変数の\textbf{同時分布（joint distribution）}の特徴に関わっています。

\begin{wideitemize}
\item 例：所得（$Y$）と教育水準（$X$）はどのように共変動するか？\pause{}

\item 同時分布のCDFは次のように定義されます：
\vspace{-0.3cm}

$$F(x,y) = Pr(X \leq x, Y \leq y)$$
\noindent すなわち、$X \leq x$ かつ $Y \leq y$ となる確率です。

\pause
\item 離散型 ($X, Y$) の場合、同時PMF $p(x,y)=Pr(X=x,Y=y)$ を定義します。

\pause
\item 連続型 ($X, Y$) の場合、同時PDF $f(x,y)=\frac{\partial^2}{\partial x \partial y}F(x,y)$ を定義します。
\end{wideitemize}
\vspace{0.5cm}
\pause{}

複数の確率変数を扱う際、個々の変数の分布を\textbf{周辺分布（marginal distribution）}と呼ぶことがあります。\pause{}

\begin{itemize}
\item 同時分布とは、例えば $p(x)=\sum_{y\in\mathbb{Y}}p(x,y)$ のような関係で結ばれています。
\end{itemize}

\end{frame}

\begin{frame}{条件付き分布}

同時分布と周辺分布を組み合わせることで、ある確率変数が与えられたときの別の確率変数の\textbf{条件付き分布（conditional distribution）}が得られます。

\begin{itemize}
\item 直感的には、条件付き分布 $Y|X=x$ とは、$X=x$ であるサブグループにおける $Y$ の分布のことです。
	
\item 条件付きPMF $p(y\mid x)=Pr(Y=y\mid X=x)=\frac{Pr(Y=y,X=x)}{Pr(X=x)}=\frac{p(y,x)}{p(x)}$\pause{}
\item 例：大学卒業を条件とした所得の分布。
\end{itemize}
\vspace{0.5cm}\pause{}

これは直ちに\textbf{ベイズの定理}を導きます：
\vspace{-0.4cm}
\begin{align*}
p(y\mid x)=p(x\mid y)\frac{p(y)}{p(x)}
\end{align*}

\end{frame}

\begin{frame}{確率変数と確率分布（続き）}
\vspace{0.2cm}
大学卒業を条件とした年収の条件付き密度関数：

\begin{center}
\includegraphics[scale=0.6]{stata12.png} 

\includegraphics[scale=0.55]{stata13.png}
\end{center}

\end{frame}

\begin{frame}{独立性}

\begin{wideitemize} 
	
\item
このコースで重要な概念となるのが\textbf{独立性（independence）}です。

\item
直感的には、独立性とは、$X$ の値を知っても $Y$ の値について何も分からないということです。

\item
形式的には、すべての $(y, x)$ に対して条件付き分布が周辺分布と一致する場合、$X$ と $Y$ は独立である（$X \indep Y$）と言います：
$$ p(y\mid x)=p(y)$$

\pause
\item 
例：$D$ がランダムに割り当てられた処置であるなら、$D \indep (Y(1),Y(0))$ です。
\pause
\begin{itemize}
	\item 
	理解度チェック：これは $D \indep Y$ を意味しますか？
\end{itemize}

\pause 

\item
\textbf{条件付き独立性}も同様に定義されます。$W$ が与えられた下で $X$ を知っても $Y$ についての情報が増えない場合、$Y \indep X \mid W$ と書きます：
$$ p(y\mid x,w)=p(y|w)$$

\end{wideitemize}
\end{frame}

\begin{frame}{多変量正規分布}

\vspace{0.2cm}
重要な多変量分布：多変量正規分布 $\mathbf{X}\in\mathbb{R}^K$

\begin{itemize}
\item 平均ベクトル $\boldsymbol{\mu}\in\mathbb{R}^K$ と、正定値な分散共分散行列 $\boldsymbol{\Sigma}\in\mathbb{R}^K\times\mathbb{R}^K$ によってパラメータ化されます。
\vspace{0.1cm}
\item 注意：原則として、ベクトルや行列は\textbf{太字}で表記します。
\end{itemize}

\vspace{0.4cm}\pause{}

多くの有用な性質がありますが、ここではいくつか紹介します。$(X,Y)^\prime$ が多変量正規分布に従うとき：

\begin{itemize}
\item $X$ と $Y$ の周辺分布は正規分布になります。
\vspace{0.1cm}
\item $X\mid Y$ および $Y\mid X$ の条件付き分布も正規分布になります。
\vspace{0.1cm}
\item 任意の線形結合 $aX+bY+c$ も正規分布に従います。
\end{itemize}

\end{frame}

\begin{frame}{2変量正規分布の密度関数}
\begin{center}
\includegraphics[scale=1.5]{binormal.jpg}
\end{center}
\end{frame}

\begin{frame}{アウトライン}
	
	\textcolor{red!75!green!50!blue!25!gray}{1. 確率変数と確率分布}$\checkmark$
	\vspace{0.8cm}
	
	2. 平均と分散 
	\vspace{0.8cm}
	
	\textcolor{red!75!green!50!blue!25!gray}{3. 実験における識別}
	\vspace{0.8cm}
	
	\textcolor{red!75!green!50!blue!25!gray}{4. 無作為抽出と標本平均}
	
	\vspace{0.8cm}
	\textcolor{red!75!green!50!blue!25!gray}{5. 仮説検定と推論}
	
\end{frame}

\begin{frame}{平均}

\begin{wideitemize}

\item
私たちはしばしば、経済的な確率変数の平均（例：世帯所得）に興味を持ちます。

\item\pause{}
$X$ の\textbf{平均（mean）}/\textbf{期待値（expectation）}とは、確率で重み付けされた代表的な値のことです。

\begin{itemize}
\item 離散型確率変数の場合：
\vspace{-0.1cm}
\begin{align*}
E[X]=\sum_{x\in\mathbb{X}}p(x)x=x_1Pr(X=x_1)+\dots + x_KPr(X=x_K)
\end{align*}\pause{}
\vspace{-0.2cm}
解釈：抽出を何度も繰り返したときの $X$ の長期的な平均値。\pause{}

\vspace{0.5cm}
\item 連続型確率変数の場合：$E[X]=\int_{x\in\mathbb{X}}f(x)xdx$\pause \\ 注意：$p(x)$ が極端な値に対して高い確率を持つ場合、期待値が存在しないこともあります。
\end{itemize}
\vspace{0.1cm}\pause

\item \textbf{重要な事実：} 期待値演算子は\textbf{線形}です。定数 $a, b$ に対して $E[a+bX]=a+bE[X]$。
\vspace{0.1cm}

\begin{itemize}
\item これは上記の定義から容易に証明できます。
\end{itemize}
\end{wideitemize}

\end{frame}

\begin{frame}{平均の計算：簡単な例}
	\begin{wideitemize}
		
		\item
		$X$ を公正なサイコロの目とします。$E[X]$ はいくらですか？
		
		\smallskip
		\pause
		\item
		定義より、
		$$E[X] = Pr(X=1) \times 1 + Pr(X=2) \times 2 + .... + Pr(X=6) \times 6$$
		
		\pause
		\item
\smallskip
		公正なサイコロなら、$Pr(X=1) =...=Pr(X=6) = \frac{1}{6}$ です。
		
		\pause 
		\item 
		これを代入すると：
		
		$$E[X] = \frac{1}{6} (1+...+6) = 3.5$$
		
	\end{wideitemize}	
\end{frame}

\begin{frame}{分散}
\vspace{0.2cm}
\textbf{分散（Variance）}は、分布の広がりの二乗を測定します：

\begin{itemize}
\item $Var(X)=E\left[(X-E[X])^2\right]\pause{}=E\left[X^2\right]-E\left[X\right]^2$ （なぜこうなるか分かりますか？）
\vspace{0.1cm}\pause{}
\item $X$ の\textbf{標準偏差（standard deviation）}は $Std(X)=\sqrt{Var(X)}$ です。これは平均からの「典型的な」ズレを表します。
\end{itemize}
\vspace{0.2cm}
\pause{}

線形変換の分散：
\begin{align*}
Var(a+bX)=&\pause{}E[(a+bX-E[a+bX])^2]\pause{} \\
 = &E[(a+bX-a-bE[X])^2]\pause{} \\
= & b^2 E[(X-E[X])^2]\pause{} \\
 =& b^2Var(X)
\end{align*}

\pause{}
これは $Std(a + bX) = |b| \cdot Std(X)$ であることを意味します。
	\begin{itemize}
		\item 
		直感的には、所得をドル単位からセント単位に変えると、標準偏差は100倍になります。
	\end{itemize}

\end{frame}

\begin{frame}{共分散}
\vspace{0.2cm}
\textbf{共分散（Covariance）}は、2つの変数の間の線形な関連性を測定します：

\begin{itemize}
\item $Cov(X,Y)=E\left[(X-E[X])(Y-E[Y])\right]=E[XY]-E[X]E[Y]$
\vspace{0.1cm}\pause{}
\item $Cov(X,X)=E\left[(X-E[X])^2\right]=Var(X)$ となることに注目してください。
\vspace{0.1cm}\pause{}
\item $X$ と $Y$ の\textbf{相関（correlation）}は $Corr(X,Y)=\frac{Cov(X,Y)}{Std(X)Std(Y)}$ です。
\end{itemize}
\vspace{0.3cm}
\pause{}

$Cov(X,Y)>0$ は、$X$ が平均より高いときに $Y$ も平均より高くなる傾向があることを意味します（逆も同様）。\pause{}

\begin{itemize}
\item $Corr(X,Y)$ は、線形な関連性を表す\textbf{単位のない}（標準化された）尺度です。
\vspace{0.1cm}\pause{}
\item $X$ と $Y$ が独立なら、$Cov(X,Y)=Corr(X,Y)=0$ です。
\vspace{0.1cm}\pause{}
\item しかし、逆は必ずしも成り立ちません！独立性は相関ゼロよりも強い概念です。
\end{itemize}

\end{frame}

\begin{frame}{相関の例}

\begin{center}
\includegraphics[scale=0.35]{corr2.png} 
\end{center}

\end{frame}

\begin{frame}{相関の例}

\begin{center}
\includegraphics[scale=0.35]{corr3.png} 
\end{center}

\end{frame}

\begin{frame}{相関の例}

\begin{center}
\includegraphics[scale=0.35]{corr4.png} 
\end{center}

\end{frame}

\begin{frame}{線形結合の平均と分散}
\begin{wideitemize}

\item
期待値は線形です：$E[aX+bY+c]=aE[X]+bE[Y]+c$ 

\pause

\item 分散は「2次的」です：$Var(aX+bY+c)=a^2Var(X)+2abCov(X,Y)+b^2Var(Y)$

\pause
\item
共分散は線形です：$Cov(aX+c,bY+d)=abCov(X,Y)$ \\ および $Cov(X+Z,Y)=Cov(X,Y)+Cov(Z,Y)$

\end{wideitemize}

\end{frame}

\begin{frame}{条件付き期待値}
\vspace{0.2cm}
経済学では、特に\textbf{条件付き期待値（conditional expectations）}に興味があります。

\begin{itemize}
\item $X=x$ のときの $Y$ の平均は何ですか？（例：ブラウン大学に進学した人の平均所得は？）
\pause 
	
\item 条件付き期待値関数 (CEF)：\\
$E[Y\mid X=x]=\sum_{y\in\mathbb{Y}} y p(y\mid x)$ （離散型の場合）\\
$E[Y\mid X=x]=\int_{y\in\mathbb{Y}} y f(y\mid x) dy$ （連続型の場合）

\vspace{0.1cm}\pause{}
\item $X$ で評価された\textbf{ランダムな}CEFを $E[Y\mid X]$ と書くことがあります。
\vspace{0.1cm}\pause{}
\item すべての $x$ に対して $E[Y\mid X=x]=E[Y]$ であるとき、$Y$ は $X$ について\textbf{平均独立（mean independent）}であると言います。
\vspace{0.1cm}\pause{}
\item $X$ で条件付けると、$X$ の関数は定数として扱えます：例：$E[f(X)+g(X)Y|X=x]=f(x)+g(x)E[Y|X=x]$。
\end{itemize}

\end{frame}

\begin{frame}{条件付き期待値の例}
	\vspace{0.2cm}
	教育年数を条件とした年収（対数）の期待値関数 (CEF)
	
	\begin{center}
		\includegraphics[scale=1.3]{stata2pre.png}
	\end{center}
	
\end{frame}

\begin{frame}{反復期待値の法則 (LIE)}
非常に重要な結果：\textbf{反復期待値の法則（Law of Iterated Expectations, LIE）}
\vspace{0.2cm}

例から始めましょう。米国の成人の平均身長を計算したいとします。LIEによれば、以下の手順で計算できます： \vspace{0.2cm}

1) 男性の平均身長を計算する。 \\ \vspace{0.2cm}

2) 女性の平均身長を計算する。 \\ \vspace{0.2cm}

3) 男女それぞれの平均身長を（男女比で重み付けして）平均する。

\vspace{0.2cm}
\pause
数学的には：

\begin{align*}
E[身長] &= P(女) E[身長 | 女] + P(男) E[身長|男] \\ \pause{}
& =  E[ E[身長 | 性別 ]  ]	
\end{align*}
 
\end{frame}

\begin{frame}{反復期待値の法則 (LIE)}
	\textbf{反復期待値の法則（LIE）} の形式的な表現は：
	\begin{align*}
		E[Y]=E[E[Y\mid X]]\pause{}
	\end{align*}
	左辺の期待値は $p(y)$ を用い、右辺の外側の期待値は $p(x)$ を、内側の期待値は $p(y\mid x)$ を用いることに注意してください。
	
\end{frame}

\begin{frame}{平均独立と無相関}
\vspace{0.2cm}
LIEを用いると、平均独立が無相関を意味することを示せます：
\begin{align*}
Corr(X,Y)&\propto E[(X-E[X])(Y-E[Y])]\\ \pause{}
&=E[E[(X-E[X])(Y-E[Y])\mid X]]\\ \pause{}
&= E[(X-E[X])E[Y-E[Y]\mid X]]\\\pause{}
&= E[(X-E[X])(E[Y\mid X]-E[Y])]\\\pause{}
&= 0\text{ （$E[Y\mid X]=E[Y]$ のとき）}
\end{align*}
\pause{}
\vspace{0.3cm}
各ステップがどのように導出されたか確認してください。
\vspace{0.5cm}
\pause{}

逆は成り立ちません：無相関な変数が平均従属であることもあります。
\pause{}

\begin{itemize}
\item もちろん、独立 $\implies$ 平均独立 です（が、逆は必ずしも成り立ちません）。
\end{itemize}

\end{frame}

\begin{frame}{ベクトル・行列表記についての補足}

\vspace{0.2cm}
確率ベクトル $\mathbf{X}=[X_1,\dots,X_N]^\prime$ を扱うと便利なことが多いです。

\begin{itemize}
\item このコースでは、これらは常に「列ベクトル」とし、太字で表記します。
\vspace{0.1cm}
\item 行列も太字で表記します（$K$ 列 $N$ 行など）。
\end{itemize}\vspace{0.4cm}\pause{}

ベクトル・行列の標準的な演算（転置、逆行列など）については、\uline{The Matrix Cookbook} が参考になります。

\begin{itemize}
\item \url{www.math.uwaterloo.ca/~hwolkowi/matrixcookbook.pdf}
\end{itemize}\vspace{0.4cm}\pause{}

期待値は各要素ごとに適用されます：例：$E\left[\begin{bmatrix}X_{11}& X_{12}\\X_{21} & X_{22}\end{bmatrix}\right]=\begin{bmatrix}E[X_{11}]& E[X_{12}]\\E[X_{21}] & E[X_{22}]\end{bmatrix}$ 
\begin{itemize}
\item 分散は $Var\left(\begin{bmatrix}X_{1}\\X_{2}\end{bmatrix}\right)=\begin{bmatrix}Var(X_{1}) & Cov(X_{1},X_{2})\\Cov(X_{1},X_{2}) & Var(X_{2})\end{bmatrix}$、共分散は $Cov\left(\begin{bmatrix}X_{1}\\X_{2}\end{bmatrix},\begin{bmatrix}Y_{1}\\Y_{2}\end{bmatrix}\right)=\begin{bmatrix}Cov(X_{1},Y_{1}) & Cov(X_{1},Y_{2})\\Cov(X_{2},Y_{1}) & Cov(X_{2},Y_{2})\end{bmatrix}$ のように定義されます。
\end{itemize}

\end{frame}

\begin{frame}{アウトライン}
	
	\textcolor{red!75!green!50!blue!25!gray}{1. 確率変数と確率分布}$\checkmark$
	\vspace{0.8cm}
	
	\textcolor{red!75!green!50!blue!25!gray}{2. 平均と分散} $\checkmark$
	\vspace{0.8cm}
	
	3. 実験における識別 
	\vspace{0.8cm}
	
	\textcolor{red!75!green!50!blue!25!gray}{4. 無作為抽出と標本平均}
	
	\vspace{0.8cm}
	\textcolor{red!75!green!50!blue!25!gray}{5. 仮説検定と推論}
	
\end{frame}

\begin{frame}{実験における識別}
\begin{wideitemize}
\item
これまでに学んだ理論を使えば、なぜ実験が（少なくとも識別の観点から）「うまくいく」のかを数学的に示すことができます。
\pause
\item \textbf{潜在的結果（potential outcomes）} の枠組みを思い出しましょう：

\begin{itemize}
	\item
	$Y_i(1), Y_i(0)$ は、個人 $i$ が処置を受けた場合と受けなかった場合の結果です。\smallskip
	
	\item
	これらを用いて観測される結果をモデル化します：$Y_i = D_i Y_i(1) + (1-D_i) Y_i(0)$ 
	
\end{itemize}

\pause
\item 
平均処置効果（ATE）に興味があるとします：
$$ATE = E[Y_i(1) - Y_i(0) ]$$

\pause
\item
各個人にコイン投げで $D_i$ を割り当てるとします。これは $D_i \indep (Y_i(1),Y_i(0))$ を意味します。なぜでしょうか？
\end{wideitemize}		
\end{frame}

\begin{frame}{ランダム化の利用}
\vspace{0.2cm}
\begin{wideitemize}
\item
実験により、$D_i \indep (Y_i(1),Y_i(0))$ が成り立っています。

\item
$\underbrace{E[Y_i | D_i = 1]}_{\text{処置群の母平均}}$ はいくらでしょうか？
\vspace{.3cm}
\pause
$$	E[Y_i|D_i=1] = E[Y_i(1) | D_i = 1] = E[Y_i(1)]$$
\noindent 最初の等号は潜在的結果モデルに基づき、2番目の等号は（平均）独立性に基づきます。

\pause
\item
同様に、$E[Y_i | D_i = 0] =  E[Y_i(0) | D_i = 0] = E[Y_i(0)]$ となります。

\pause
\item
これらの結果を組み合わせると：
\vspace{-0.2cm}

$$\underbrace{E[Y_i | D_i = 1]}_{\text{処置群の母平均}} - \underbrace{E[Y_i | D_i = 0]}_{\text{対照群の母平均}} = \underbrace{E[Y_i(1) - Y_i(0)]}_{\text{平均処置効果}} = \tau$$

したがって、実験における処置群と対照群の母平均の差によって ATE が識別されます！
\end{wideitemize}
\end{frame}

\begin{frame}{条件付き非交絡性（Unconfoundedness）の下での識別}
	\begin{wideitemize}
		\item
		ここで、$D_i \indep (Y_i(1), Y_i(0)) | \mathbf{X}_i$ であると仮定します。ただし $\mathbf{X}_i$ は観察可能な特徴のベクトルです。
		
		\item
		これは条件付き非交絡性、あるいは\textbf{観測される変数に基づく選択（selection on observables）}と呼ばれます。
		
		\pause
		\item
		直感的には、条件付き非交絡性とは、同じ $\mathbf{X}_i$ の値を持つ人々の間では、実質的に実験が行われているのと同じ状況であることを意味します。
\begin{itemize}
\item これは（無条件の）独立性から導かれますが、より\textbf{弱い}条件です。$\mathbf{X}_i$ を通じた非ランダム性を許容するからです。
\end{itemize}
		
		\pause
		\item 
		なぜ条件付き非交絡性を信じられるのでしょうか？
		\begin{itemize}
			\item 
			層別実験（Stratified experiment）：同じ $\mathbf{X}_i$ を持つ人々の間でランダム化を行う場合。
			
			\medskip
			\pause
			\item
			「擬似実験（Quasi-experiment）」/「自然実験」：同じ $\mathbf{X}_i$ を持つ人々の間では、$D_i$ が実質的にランダムに決まっていると考えられる場合。
		\end{itemize}
	\end{wideitemize}
\end{frame}

\begin{frame}{例：猛暑日とテストの点数}
	
\begin{wideitemize}
\item
\href{https://www.nature.com/articles/s41562-020-00959-9}{\uline{Park et al (2021)}} は、学期中の猛暑日（$D_i$）がテストの点数（$Y_i$）に与える影響を研究しました。\smallskip
\begin{itemize}
	\item 
	彼らの $D_i$ は2値変数ではありませんが、例えば $D_i = 1[猛暑日 > 10日]$ のように2値化して考えることもできます。
\end{itemize}

\item
なぜ $D_i \not\indep (Y_i(1),Y_i(0))$ だと考えられるのでしょうか？
\pause
	\begin{itemize}
		\item 
		気候の異なる場所に住む人々は、他の側面（経済状況など）でも異なっている可能性があるからです。
	\end{itemize}

\pause
\item
Park らは、$X_i$ をその場所の過去の気象データとしたとき、$D_i \indep (Y_i(1),Y_i(0)) | X_i$ であると主張しています。

	\begin{itemize}
		\item 
		過去の気象パターンを条件とすれば、特定の年の暑さは実質的にランダムであるという考えです。
	\end{itemize}

\item\pause{}
これは妥当な仮定だと思いますか？

\end{wideitemize}	
	
\end{frame}

\begin{frame}{Park et al. (2021) 要約}
	\centering
	\includegraphics[width =0.6\linewidth]{park-heat}
\end{frame}

\begin{frame}{例：大学の選択度（Selectivity）の収益}
\begin{wideitemize}
\item
\href{https://academic.oup.com/qje/article-abstract/117/4/1491/1876022}{\uline{Dale and Krueger (2002)}} は、私たちの例と似た問いを研究しました：\\
選択度の高い大学に通うことは、所得にどのような影響を与えるか？

\item 
明らかに $D_i \not\indep (Y_i(1),Y_i(0))$ です。なぜなら、選択度の高い大学に通う学生は、もともとの学力も高い傾向があるからです。

\pause
\item
Dale と Krueger は、$\mathbf{X}_i$ を「その学生が出願し、合格した大学の集合」としたとき、$D_i \indep (Y_i(1),Y_i(0)) | \mathbf{X}_i$ であると主張しました。

\item
本質的に、彼らは「どの大学に出願し、どこに合格したか」さえ分かれば、その中からどこに進学するかは実質的にランダムであると議論しています。

\pause
\item
これを信じますか？この仮定が崩れるとしたら、どのような理由が考えられますか？
\pause
	\begin{itemize}
		\item 
		合格した大学の中からどこを選ぶかは、家庭環境、意欲、将来の計画などの違いを反映しているかもしれません。
	\end{itemize}
\end{wideitemize}
\end{frame}

\begin{frame}{Dale and Krueger (2002)}
	\centering
	\includegraphics[width = 0.9\linewidth]{dale-krueger}
\end{frame}

\begin{frame}{条件付き非交絡性の利用}
\vspace{0.2cm}
	\begin{wideitemize}
		\item
		$D_i \indep (Y_i(1), Y_i(0)) | \mathbf{X}_i$ であるとします。
		
		\item
		実験の場合と同様に、すべての $\mathbf{x}$ に対して次が成り立ちます：
		$$E[Y_i | D_i = 1,\mathbf{X}_i= \mathbf{x}] = E[Y_i(1) | \mathbf{X}_i=\mathbf{x}]$$
		
		および
		
		$$E[Y_i | D_i = 0,\mathbf{X}_i = \mathbf{x}] = E[Y_i(0) | \mathbf{X}_i =\mathbf{x}]$$
		
		\pause
		\item
		これは次を意味します：
		$$\small \underbrace{E[Y_i | D_i = 1, \mathbf{X}_i= \mathbf{x}] }_{\text{$\mathbf{X}_i= \mathbf{x}$ の処置群平均}}- \underbrace{E[Y_i | D_i = 0, \mathbf{X}_i= \mathbf{x}]}_{\text{$\mathbf{X}_i= \mathbf{x}$ の対照群平均}} = \underbrace{E[Y_i(1) - Y_i(0) | \mathbf{X}_i= \mathbf{x}]}_{\text{$\mathbf{X}_i= \mathbf{x}$ での ATE}} $$ 
		
		\pause
		\item 
		$E[Y_i(1) - Y_i(0) | \mathbf{X}_i= \mathbf{x}]$ はしばしば\textbf{条件付き平均処置効果（CATE）}と呼ばれ、$CATE(\mathbf{x})$ と書かれます。
	\end{wideitemize}	
\end{frame}

\begin{frame}{条件付き非交絡性の利用（続き）}
\vspace{0.2cm}
	\begin{wideitemize}
		\item 
		条件付き非交絡性の下で、$CATE(\mathbf{x})$ が識別されることを示しました。
		
		\item 
		では、無条件の $ATE=E[Y_i(1) - Y_i(0)]$ も識別されるでしょうか？
		
		\pause
		\item
		はい！反復期待値の法則（LIE）より：
		$$E[ \underbrace{E[Y_i(1) - Y_i(0) | \mathbf{X}_i] }_{ CATE(\mathbf{X}_i) }  ] = E[Y_i(1) - Y_i(0)]$$
		
		\pause
		
		\item
		\textbf{技術的な注意}：ここでは、すべての $\mathbf{x}$ について条件付き期待値が存在すると仮定しています。
		
		\pause
		\item
		これには、$0< Pr(D_i = 1 |\mathbf{X}_i= \mathbf{x}]) < 1$ という条件が必要です。これは\textbf{オーバーラップ（overlap）}条件と呼ばれます。
		
		\item
		直感的には、全体の ATE を知るためには、すべての $X_i$ の値において処置群と対照群の両方が存在する必要があります。
	\end{wideitemize}
\end{frame}

\begin{frame}{母平均について学ぶ}
	
	\begin{wideitemize}
		\item 実験において、平均処置効果が母平均の差として識別されることを示しました：
		
		$$E[Y_i | D_i = 1] - E[Y_i |D_i =0] = E[Y_i(1)-Y_i(0)]$$

		\item
		同様に、条件付き非交絡性の下では、CATE は条件付き母平均の差によって識別されます。
		
		\pause
		\item
		しかし、現実には母集団全体のデータは見えないため、$E[Y_i | D_i = 1]$ や $E[Y_i | D_i = 0]$ などを直接知ることはできません。
		
		\pause
		\item
		私たちは、観察されたサンプルからこれらの\textbf{推定対象（estimands）}について学ぶ必要があります。
		
		\item
		ここから\textbf{統計的推論}の出番です...
	\end{wideitemize}
	
\end{frame}

\begin{frame}{アウトライン}
	
	\textcolor{red!75!green!50!blue!25!gray}{1. 確率変数と確率分布} $\checkmark$
	\vspace{0.8cm}
	
	\textcolor{red!75!green!50!blue!25!gray}{2. 平均と分散} $\checkmark$
	\vspace{0.8cm}
	
	\textcolor{red!75!green!50!blue!25!gray}{3. 実験における識別}  $\checkmark$
	\vspace{0.8cm}
	
	4. 無作為抽出と標本平均
	
	\vspace{0.8cm}
	\textcolor{red!75!green!50!blue!25!gray}{5. 仮説検定と推論}
	
\end{frame}

\begin{frame}{サンプルの定義}
	\begin{wideitemize}	
	
	\item
	統計的推論を定式化するために、観察されたデータが母集団からどのように抽出されたかを特定する必要があります。
	
	\pause{}
	
	\item 基本的なケース：サイズ $N$ の\textbf{独立で同一な分布（iid）}に従う\textbf{代表的な（representative）}サンプルを観察するとします：例：$\mathbf{Y}=[Y_{1},Y_{2},\dots,Y_{N}]^\prime$
	
	\pause{}\smallskip
\begin{itemize}
	\item \textbf{独立（Independent）}：すべての $i\neq j$ について $Y_i$ は $Y_j$ と独立。
	
\smallskip
	\pause{}
	\item \textbf{同一な分布（Identically distributed）}：すべての $i, j$ について $Y_i$ と $Y_j$ は同じ分布に従う。

\pause{}
\smallskip
	\item
	\textbf{代表的（Representative）}：$Y_i$ の分布が、関心のある母集団の分布と同じである。
	
\end{itemize}		
	\pause	
	\item	
	iid かつ代表的なデータは分析が容易な基準となりますが、現実には必ずしも成り立たないことを意識する必要があります。
		\begin{itemize}
			\pause
\smallskip
			\item
			同じ世帯の人々を一緒に抽出する場合、独立ではありません！
		
\smallskip	
			\pause
			\item
			州ごとに層化抽出を行う場合、同一分布ではありません。
			\pause

\smallskip
			\item 
			デューイ対トルーマンの例では、代表的ではありません！
		\end{itemize}
	\end{wideitemize}

\end{frame}

\begin{frame}{標本平均の期待値と分散}
\vspace{0.1cm}
母平均 $\mu=E[Y_i]$ を、サイズ $N$ の iid かつ代表的なサンプル $\mathbf{Y}$ から学びたいとします。\pause{} 
\pause{}
\begin{itemize}
\item 自然な推定量は\textbf{標本平均}です：$\hat{\mu}=\frac{1}{N}\sum_i Y_{i}$
\vspace{0.05cm}\pause{}
\item $\hat{\mu}$ はランダムなデータ $\mathbf{Y}$ の関数です。したがって、それ自体が確率変数であり、分布（「標本分布」）を持ちます。
\end{itemize}
\vspace{0.1cm}
\pause{}

これまでに学んだことを使って、$\hat{\mu}$ の期待値と分散を導出できます：
\begin{align*}
E[\hat{\mu}]&=E\left[\frac{1}{N}\sum_i Y_i\right]=\pause\frac{1}{N}\sum_i E[Y_i]=\pause\mu\\\pause{}
Var(\hat{\mu})&=Var\left(\frac{1}{N}\sum_i Y_i\right)=\pause{}\frac{1}{N^2}\sum_iVar(Y_i)=\pause{}\sigma^2/N
\end{align*}
（ただし $\sigma^2=Var(Y_i)$）
\pause{}
\vspace{-0.2cm}
\begin{itemize}
\item 式 (1) は $\hat{\mu}$ が\textbf{不偏}であることを示しています：平均値は $\mu$ です。
\pause{}
\item 式 (2) は $\hat{\mu}$ の分散がサンプルサイズ $N$ とともに減少することを示しています（\textbf{一致性}に近い概念です）。
\end{itemize}

\end{frame}

\begin{frame}{無作為抽出と標本平均}
\vspace{0.2cm}
平均所得の推定における不偏性と一致性のシミュレーション：

\begin{center}
\includegraphics[scale=0.45]{stata18.png} \includegraphics[scale=0.5]{stata19.png}
\end{center}

\end{frame}

\begin{frame}{無作為抽出のシミュレーション}
\vspace{0.2cm}
平均所得の推定における不偏性と一致性のシミュレーション：

\begin{center}
\includegraphics[scale=0.45]{stata18.png} \includegraphics[scale=0.5]{stata20.png}
\end{center}

\end{frame}

\begin{frame}{無作為抽出と標本平均}
\vspace{0.2cm}
平均所得の推定における不偏性と一致性のシミュレーション：

\begin{center}
\includegraphics[scale=0.45]{stata18.png} \includegraphics[scale=0.5]{stata21.png}
\end{center}

\end{frame}

\begin{frame}{無作為抽出と標本平均}
\vspace{0.2cm}
平均所得の推定における不偏性と一致性のシミュレーション：

\begin{center}
\includegraphics[scale=0.45]{stata18.png} \includegraphics[scale=0.5]{stata22.png}
\end{center}

\end{frame}

\begin{frame}{無作為抽出と標本平均}
\vspace{0.2cm}
平均所得の推定における不偏性と一致性のシミュレーション：

\begin{center}
\includegraphics[scale=0.45]{stata18b.png} \includegraphics[scale=0.5]{stata23.png}
\end{center}

\end{frame}

\begin{frame}{無作為抽出と標本平均}
\vspace{0.2cm}
まとめ：$\hat{\mu}=\frac{1}{N}\sum_i Y_i$ が $\mu=E[Y_i]$ の良い推定量である2つの理由：

\begin{itemize}
\item 不偏である：$E[\hat{\mu}]=\mu$
\vspace{0.1cm}
\item サンプルサイズが大きくなるにつれて分散がゼロに収束する：$\lim_{N\rightarrow\infty}Var(\hat{\mu})=0$ 
\end{itemize}
\vspace{0.5cm}
\pause{}

次章では、$\hat{\mu}$ のもう一つの優れた性質を見ます：$N$ が大きいとき、その\textbf{分布}は近似的に正規分布に従います。

\end{frame}

\begin{frame}{無作為抽出と標本平均}
\vspace{0.2cm}
条件付き期待値 $\mu(x)=E[Y_i\mid X_i=x]$ に関心がある場合、$X_i=x$ を条件とした $Y_i$ の条件付き標本平均も考えられます。
\pause{}

\begin{itemize}
\item これは $X_i$ が少ない値 $x$ を持つ離散型の場合に最も簡単です。
\vspace{0.1cm}\pause{}
\item 自然な推定量 $\hat{\mu}(x)=\frac{1}{N_x}\sum_{i:X_i=x}Y_i$ （ただし $N_x$ は $X_i=x$ である観測数）。
\end{itemize}
\vspace{0.4cm}
\pause{}

先ほどと同様の導出により：
\begin{align*}
E[\hat{\mu}(x)]&=E[Y_i\mid X_i=x]=\mu(x)\\
Var(\hat{\mu}(x))&=Var(Y_i\mid X_i=x)/N_x
\end{align*}
\pause{}
\vspace{-0.2cm}

したがって、これも不偏推定量であり、$N_x\rightarrow\infty$ となれば真の値に近づきます。
\pause{}

\begin{itemize}
\item $X_i$ が離散型でない場合や、$N_x$ が大きくならない場合はどうなるでしょうか？それは後ほど...
\end{itemize}

\end{frame}

\begin{frame}{アウトライン}
	
	\textcolor{red!75!green!50!blue!25!gray}{1. 確率変数と確率分布} $\checkmark$
	\vspace{0.8cm}
	
	\textcolor{red!75!green!50!blue!25!gray}{2. 平均と分散} $\checkmark$
	\vspace{0.8cm}
	
	\textcolor{red!75!green!50!blue!25!gray}{3. 実験における識別} $\checkmark$
	\vspace{0.8cm}
	
	\textcolor{red!75!green!50!blue!25!gray}{4. 無作為抽出と標本平均} $\checkmark$
	
	\vspace{0.8cm}
	5. 仮説検定と推論 
	
\end{frame}

\begin{frame}{仮説検定 -- 導入}
\begin{wideitemize}

\item
$N$ が大きくなれば、標本平均 $\hat\mu$ は母平均 $\mu$ に近づくことを示しました。

\item
しかし、「近い」とはどういう意味でしょうか？

\item
データの所得の標本平均が50,000ドルだったとき、母平均が55,000ドルであると考えるのは妥当でしょうか？70,000ドルならどうでしょうか？

\pause
\item
\textbf{仮説検定（Hypothesis testing）}は、この「近い」という概念を定式化するのに役立ちます。

\pause
\item
それは、もし真の値が55,000ドルや70,000ドルだった場合に、標本平均が50,000ドルになる確率がどの程度あるかを教えてくれます。

\end{wideitemize}
\end{frame}

\begin{frame}{仮説検定の概要}

\begin{enumerate}
	\item 
	母平均が特定の値であるという\textbf{帰無仮説}を指定します：$H_0: \mu = \mu_0$。
\smallskip
		\begin{itemize}
			\item 
			例：母平均が55,000ドルなら $H_0: \mu = 55,000$。
			\end{itemize}
	
	\medskip
	\pause
	\item
	帰無仮説が正しいとした場合に、$\hat\mu$ が $\mu_0$ から少なくともこれほど離れた値として観察される確率を計算します。これを \textbf{p値（p-value）} と呼びます。
	\medskip
	\pause
	\item
	p値が小さい場合、帰無仮説を\textbf{棄却（reject）}します。すなわち、帰無仮説が正しいとしたら、これほど $\mu_0$ から離れた $\hat\mu$ が観察されるのは考えにくいということです。
\smallskip
		\begin{itemize}
		\pause 
		\item
		一般的なしきい値は $\alpha=0.05$ です。
	\end{itemize}
	
	\vspace{.1cm}
	\pause
	\item
	このようにして棄却できないすべての $\mu_0$ の値を集めて、\textbf{信頼区間（confidence interval, CI）}を形成します。
		\begin{itemize}
			\pause
			\item
			$\alpha = 0.05$ の場合、構築された信頼区間は、データの実現値の 95\% において真の値 $\mu$ を含みます。
		\end{itemize}
\end{enumerate}
	
\end{frame}

\begin{frame}{標本平均 $\hat\mu$ が正規分布に従う場合の仮説検定}
\begin{itemize}
	\item 
	特別なケースとして、$\hat\mu \sim \mathrm{N}(\mu, \sigma^2/N)$ かつ $\sigma^2$ が既知である場合を考えましょう。
	\medskip
	\pause
	\item
	なぜこのケースを考えるのでしょうか？
\smallskip
		\begin{itemize}
			\item 
			$Y_i \overset{iid}{\sim} \mathrm{N}(\mu,\sigma^2)$ なら、これが $\hat\mu$ の正確な分布になります。\pause
			\smallskip
			\item
			次章で示しますが、$Y_i$ が正規分布でなくても、$N$ が大きければ $\hat\mu$ は近似的に正規分布に従います。
\smallskip\pause{}
\item また、$N$ が大きければ $\sigma^2$ を非常に精度よく推定できることも示します。
		\end{itemize}

	\pause
	\item
	帰無仮説 $H_0:\mu=\mu_0$ （例：平均所得は55,000ドル）を検定したいとします。
	
	\item
	$\hat{t} = \dfrac{\hat\mu - \mu_0 }{ \sigma /\sqrt{N} }$ とすると、帰無仮説の下で $\hat{t} \sim \mathrm{N}(0,1)$ となります。
	\begin{itemize}
	\pause
	\item
	$\hat{t}$ の分布は、サンプル $(Y_1,...,Y_N)$ の抽出を何度も繰り返したときの分布です。
		\end{itemize}
\end{itemize}
\end{frame}

\begin{frame}{標本平均 $\hat\mu$ が正規分布に従う場合の仮説検定（続き）}
\begin{wideitemize}
\item
帰無仮説 $H_0:\mu=\mu_0$ の下で、$\hat{t} \sim \mathrm{N}(0,1)$ であることを示しました。

\pause
\item ある $t \geq 0$ に対して、$Pr( |\hat{t}| > t )$ はいくらでしょうか？ \pause
$$ Pr( |\hat{t}| > t ) = 1 - Pr(|\hat{t}| \leq t) \pause = 1 - Pr(  -t \leq \hat{t} \leq t  ) = \pause 1 - (\Phi(t) - \Phi(-t))$$ 
\vspace{-0.8cm}
\item
帰無仮説 $H_0: \mu = \mu_0$ に対する p値を次のように定義します：
\begin{align*}
 p(\hat{t}) = 1 - (\Phi( |\hat{t} | ) - \Phi(- |\hat{t} |  )  ) \pause{} &= 1-  \left(  \Phi \left(  \frac{|\hat\mu - \mu_0|}{\sigma / \sqrt{N}} \right)  -  \Phi \left(  \frac{-|\hat\mu - \mu_0|}{\sigma / \sqrt{N}} \right) \right) 	\\
 &\pause{}= 2\left(1 -  \Phi \left(  \frac{|\hat\mu - \mu_0|}{\sigma / \sqrt{N}} \right)  \right) 
\end{align*}

\item
直感的には、p値は、帰無仮説が正しいときに少なくともこれほど大きな $|\hat{t}|$ が観察される確率です。

\end{wideitemize}	
\end{frame}

\begin{frame}{p値構築のイラスト}
	
	\begin{center}
		標準正規分布の密度関数（平均0、標準偏差1）
		\includegraphics[scale=0.3]{pvalue1.png}
	\end{center}
	
\end{frame}

\begin{frame}{p値構築のイラスト}
	
	\begin{center}
		ランダムな推定量を標準化した値（実現値）
		\includegraphics[scale=0.3]{pvalue2.png}
	\end{center}
	
\end{frame}

\begin{frame}{p値構築のイラスト}		
	\begin{center}
		p値の計算
		\includegraphics[scale=0.3]{pvalue3.png}
	\end{center}
\end{frame}

\begin{frame}{いつ帰無仮説を棄却するか？}

\begin{wideitemize}
	\item 
	p値の形式を思い出してください：
	$$1 - \left(  \Phi(|\hat{t}|) - \Phi( -|\hat{t}| )  \right) $$
	
	\item
	$\Phi(1.96) - \Phi(-1.96) \approx 0.95$ となることが知られています。したがって、$p < 0.05$ となるのは $|\hat{t}| > 1.96$ のとき、かつそのときに限られます。つまり、5\%有意水準で $|\hat{t}| > 1.96$ なら棄却します。
	
	\pause
	\item
	これは、どのような $\mu_0$ を棄却し、どのような $\mu_0$ を棄却しないことを意味するでしょうか？
	
	\pause
	\item
	次の場合に棄却しません：
	$$ |\hat{t}| \leq 1.96 \pause \implies \dfrac{ |\hat\mu - \mu_0| }{ \sigma / \sqrt{N} } \leq 1.96 \pause \implies \mu_0 \in  [\hat\mu - 1.96 \sigma / \sqrt{N} , \hat\mu + 1.96 \sigma / \sqrt{N}]$$
\vspace{-0.5cm}
	\pause
	\item
	したがって、区間 $\hat\mu \pm 1.96 \sigma / \sqrt{N}$ が 95\% 信頼区間 (CI) となります。
	\smallskip

		\begin{itemize}
			\item 
			これは、$H_0: \mu = \mu_0$ が正しいときに $Pr( \mu_0 \in CI ) = 0.95$ となる性質を持ちます。
		\end{itemize}
	
\end{wideitemize}
	
\end{frame}

\begin{frame}{有意水準と検出力}
	
	\begin{itemize}
		\item 検定の\textbf{有意水準（significance level）}（または\textbf{サイズ}）とは、帰無仮説が正しいときに誤って棄却してしまうあらかじめ設定された確率です（第I種の過誤の発生率）。
		 \pause
		 \begin{itemize}
		 	\item 
		 	例：5\%水準の検定は $p<0.05$ のときに棄却します。
		 \end{itemize}
		\pause
		\bigskip

		\item 検定の\textbf{検出力（power）}とは、帰無仮説が偽であるときに正しく棄却できる確率です（1 - 第II種の過誤の発生率）。
		
			\begin{itemize}
				\item 
				検出力は\textbf{対立仮説}の関数です。すなわち、実際には $\mu = \mu_A$ であるときに $H_0:\mu=\mu_0$ を棄却する確率です。
			\end{itemize}
	\end{itemize}

\end{frame}

\begin{frame}{p値の解釈に関する注意}

\begin{itemize}
\item 頻度論的なp値はしばしば「帰無仮説が正しい確率」と解釈されますが、これは正しいでしょうか？\pause{} いいえ！

\begin{wideitemize}
	\item p値は、帰無仮説が正しいと\textbf{仮定した}場合に、観察されたデータが得られる確率を示しています。
	
	\item つまり、p値は $P(data | H_0)$ について述べており、$P(H_0 | data)$ ではありません。

	\item ベイズの定理によれば、$P(H_0 | Data) = P(Data | H_0) * P(H_0) / P(Data)$ です。これを定式化するには、帰無仮説が正しいという\textbf{事前}の信念 $P(H_0)$ を設定する必要があります。
	
\end{wideitemize}

\pause
\bigskip
\item $p<0.05$ なら「効果がある」強い証拠であり、$p \geq 0.05$ なら「効果がない」証拠であると解釈されがちです。

	\begin{itemize}
		\item 
		しかし、0.05 というしきい値はかなり恣意的なものです。
		\vspace{0.1cm}
		\item
		p値は、帰無仮説の下で観察されたデータがどの程度尤もらしいかを示すスペクトラムとして捉えるのが適切です。\pause
		\vspace{0.1cm}
		\item
		さらに、帰無仮説が偽であっても（検出力が低ければ）p値が大きくなることもあります。
	\end{itemize}
\end{itemize}

\end{frame}

\begin{frame}
	\includegraphics[width = 0.5 \linewidth]{xkcd_sun}
\end{frame}

\begin{frame}

\begin{center}
\includegraphics[scale=0.45]{p-values0.jpg}
\end{center}

\end{frame}

\begin{frame}

\vspace{0.8cm}
\begin{center}
\includegraphics[scale=0.45]{p-values.jpg}
\end{center}

\end{frame}

\begin{frame}{非正規データの場合は？}
\vspace{0.2cm}
\begin{itemize}
\item 分散が既知の正規確率変数の平均について、仮説検定や推論の方法が分かりました。...だから何だと言うのでしょうか？
\pause{}
\smallskip

\begin{itemize}
\item ほとんどの変数は正規分布に従いません！\pause{}
\vspace{0.1cm}
\item たとえ正規分布だとしても、なぜ分散が既知だと言えるのでしょうか？！\pause{}
\vspace{0.1cm}
\item 私は皆さんの時間を無駄にさせてしまったのでしょうか！？ \pause{} いいえ。\pause{}
\end{itemize}
\vspace{0.4cm}

\item 次に、強力な\textbf{漸近（asymptotic）}理論を概観します。これにより、サンプルが「十分に大きい」場合には、$Y_i$ が正規分布に従わなくても、同様の推論ツールを適用できるようになります。
\end{itemize}

\end{frame}

\end{document}
